\section{Discussion}\label{sec:discussion}

\subsection{Extensions}

\paragraph{Heterogeneity in means} Letting the mean of a random parameter depend on site characteristics, $\beta_{ij}=\beta_j+\bm w_i^\top\bm\phi_j+\sigma_j z_{ij}$, adds interaction terms $x_{itj}w_{ik}$ that enter $\psi_{it}$ linearly. They can receive spike-and-slab priors with the heredity constraint $\phi_{jk}\neq0$ only if $\delta_j=1$, and are updated exactly as the $\beta_j$ are.

\paragraph{Correlated random parameters} Writing $\varepsilon_{it}=\x_{it}^\top L\bm z_i$ with $L$ lower triangular, each element $L_{jk}$ enters the linear predictor linearly given $\bm z_i$, so spike-and-slab priors on the off-diagonal elements select correlations in the manner of \citet{fruhwirth2008bayesian}. The collapsed step of Proposition~\ref{prop:marginal} must then integrate over the correlated components jointly. \tbd{Work out whether a rank-one structure keeps this tractable.}

\paragraph{Heterogeneity in variances} A site-dependent standard deviation such as
\[
\sigma_{ij}=\sigma_j\exp(\bm w_i^\top\bm\theta_j)
\]
is no longer linear in $\bm\theta_j$ and requires a Metropolis step, but the selection logic is unchanged.

\paragraph{Other outcomes} Injury severity models with random parameters are logit or ordered models. P\'olya-Gamma augmentation applies directly to binary and multinomial logit likelihoods \citep{polson2013bayesian}, so the three-state prior and the collapsed update carry over.

\paragraph{Scale} For very large networks, the expectation-maximization approach to spike-and-slab selection \citep{rockova2014emvs, rockova2018spike} can locate high-posterior specifications quickly, with the sampler used to quantify uncertainty around them.

\subsection{Limitations}

The random parameters are normal. Heavier-tailed or multimodal heterogeneity, for example finite mixtures \citep{buddhavarapu2016modeling}, would change the variance function in Proposition~\ref{prop:cross} and the integral in Proposition~\ref{prop:marginal}. One-at-a-time indicator updates can mix slowly when candidates are strongly correlated. Most importantly, the method cannot create information that the data do not contain. In cross-sectional data, a selected random parameter remains compatible with covariate-dependent dispersion, and the posterior state probabilities should be read in that light. The method makes this uncertainty visible rather than resolving it.
